% Präambel
\documentclass[12pt,a4paper,oneside, 
liststotoc, 					% Tabellen- und Abbildungsverzeichnis ins Inhaltsverzeichnis
bibtotoc,						% Literaturverzeichnis ins Inhaltsverzeichnis aufnehmen
titlepage, 						% Titlepage-Umgebung statt \maketitle
headsepline, 					% horizontale Linie unter Kolumnentitel
%abstracton,					% Überschrift beim Abstract einschalten, Abstract muss dazu in {abstract}-Umgebung stehen
%DIV11,							% auskommentieren, um den Seitenspiegel zu vergrößern
BCOR6mm,						% Bindekorrektur, die den Seitenspiegel um 6mm nach rechts verschiebt,
]{scrreprt}		

% Laden verschiedener Packages
\usepackage{ucs} 				% Dokument in utf8-Codierung schreiben und speichern
\usepackage{url}
\usepackage[utf8x]{inputenc} 	% ermöglicht die direkte Eingabe von Umlauten
\usepackage[ngerman]{babel} 	% deutsche Trennungsregeln und Übersetzung der festcodierten Überschriften
\usepackage[T1]{fontenc} 		% Ausgabe aller zeichen in einer T1-Codierung (wichtig für die Ausgabe von Umlauten!)
\usepackage{graphicx}  			% Einbinden von Grafiken erlauben
\usepackage{amsmath}
\usepackage{csquotes}           % Zusatzpaket für Zitierungen (Nur optisch, nicht logisch)
%\usepackage{amsfonts}
%\usepackage{amssymb}
\usepackage{mathpazo} 			% Einstellung der verwendeten Schriftarten
\usepackage{textcomp} 			% zum Einsatz von Eurozeichen u. a. Symbolen
\usepackage{listings}			% Datstellung von Quellcode mit den Umgebungen {lstlisting}, \lstinline und \lstinputlisting
\usepackage[dvipsnames]{xcolor} % einfache Verwendung von Farben in nahezu allen Farbmodellen
\usepackage[intoc]{nomencl} 	% zur Erstellung des Abkürzungsberzeichnisses
\usepackage{fancyhdr}			% Zusatzpaket zur Gestaltung von Fuß und Kopfzeilen
\usepackage{epigraph}
\usepackage{svg}
\usepackage{tkz-kiviat}

\usepackage{xargs}
\usepackage{setspace}
\usepackage[smaller]{acronym}

\usepackage[]{algorithm2e} 		% Für die Verwendung von Pseudocode
\usepackage[]{algorithmicx}		% Anderes Pseudocode Paket zum testen

\usepackage[colorinlistoftodos,prependcaption,textsize=tiny]{todonotes}

\usepackage{hyperref}

%Einstellungen, Farben, Commands, 
\newcommandx{\unsure}[2][1=]{\todo[linecolor=red,backgroundcolor=red!25,bordercolor=red,#1]{#2}}
\newcommandx{\change}[2][1=]{\todo[linecolor=blue,backgroundcolor=blue!25,bordercolor=blue,#1]{#2}}
\newcommandx{\info}[2][1=]{\todo[linecolor=OliveGreen,backgroundcolor=OliveGreen!25,bordercolor=OliveGreen,#1]{#2}}
\newcommandx{\improvement}[2][1=]{\todo[linecolor=Plum,backgroundcolor=Plum!25,bordercolor=Plum,#1]{#2}}
\newcommandx{\thiswillnotshow}[2][1=]{\todo[disable,#1]{#2}}

% Default fixed font does not support bold face
\DeclareFixedFont{\ttb}{T1}{txtt}{bx}{n}{12} % for bold
\DeclareFixedFont{\ttm}{T1}{txtt}{m}{n}{12}  % for normal

%------------Colors ----------------
\definecolor{deepblue}{rgb}{0,0,0.5}
\definecolor{deepred}{rgb}{0.6,0,0}
\definecolor{deepgreen}{rgb}{0,0.5,0}

\pgfkeys{
	/kiviatgrad/simplify label/.code={
		\ifx\nv\undefined\else
		\pgfmathparse{Mod(\nv,5)}
		\ifdim\pgfmathresult pt>0pt
		\tikzset{opacity=0}
		\fi
		\fi
	}
}	
% ------------------Setting up environment for codes ---------------
% Python style for highlighting
\newcommand\pythonstyle{\lstset{
		language=Python,
		basicstyle=\ttm,
		morekeywords={access,and,break,class,continue,def,del,elif,else,except,exec,finally,for,from,global,if,import,in,is,lambda,not,or,pass,print,raise,return,try,while},
		morekeywords=[2]{abs,all,any,basestring,bin,bool,bytearray,callable,chr,classmethod,cmp,compile,complex,delattr,dict,dir,divmod,enumerate,eval,execfile,file,filter,float,format,frozenset,getattr,globals,hasattr,hash,help,hex,id,input,int,isinstance,issubclass,iter,len,list,locals,long,map,max,memoryview,min,next,object,oct,open,ord,pow,property,range,raw_input,reduce,reload,repr,reversed,round,set,setattr,slice,sorted,staticmethod,str,sum,super,tuple,type,unichr,unicode,vars,xrange,zip,apply,buffer,coerce,intern},%
		morecomment=[l]\#,%
		sensitive=true,%
		%otherkeywords={self,if,print,while,return,def},             % Add keywords here
		keywordstyle=\ttb\color{deepblue},
		emph={MyClass,__init__},          % Custom highlighting
		emphstyle=\ttb\color{deepred},    % Custom highlighting style
		stringstyle=\color{deepgreen},
		frame=tb,                         % Any extra options here
		showstringspaces=false            % 
}}

% Python environment
\lstnewenvironment{python}[1][]
{
	\pythonstyle
	\lstset{#1}
}
{}

% Python for external files
\newcommand\pythonexternal[2][]{{
		\pythonstyle
		\lstinputlisting[#1]{#2}}}


\usetikzlibrary{arrows} % define style of tkiz kiviat

% Abkürzungen
\newcommand{\ua}{\mbox{u.\,a.\ }}
\newcommand{\zB}{\mbox{z.\,B.\ }}
\newcommand{\bs}{$\backslash$}

\renewcommand{\nomname}{Abkürzungsverzeichnis}

% -------------------------------------------------------------------------------------------
% Definition der Kopf- und Fußzeilen
\lhead{}								% Kopf links
\chead{}								% Kopf mitte
\rhead{\sffamily{\leftmark}}				% Kopf rechts
\lfoot{}								% Fuß links
\cfoot{\sffamily{\thepage}}				% Fuß mitte
\rfoot{\sffamily{\autor}}				% Fuß rechts
\renewcommand{\headrulewidth}{0.4pt}	% Liniendicke Kopf
\renewcommand{\footrulewidth}{0.4pt}	% Liniendicke Fuß


\makenomenclature							% Abkürzungsverzeichnis erstellen
%\input{Inhalt/abkuerzungen}					% Datei mit Abkürzungen laden

% Definition of colors
\definecolor{lightblue}{rgb}{0.910,0.933,0.970}
\definecolor{lightred}{RGB}{247,238,232}
\definecolor{monochromeLightblue}{RGB}{165,188,222}
\definecolor{monochromeLightred}{RGB}{222,188,165}
\definecolor{kiviatOne}{RGB}{137,193,30}
\definecolor{kiviatTwo}{RGB}{20,128,120}
\definecolor{kiviatThree}{RGB}{208,108,32}
\definecolor{kiviatFour}{RGB}{168,26,104}

\newcommand\ColorBox[1]{\textcolor{#1}{\rule{2.5ex}{2.5ex}}}

% ----------------------------------- Links Styling -----------------------------------------
\hypersetup{
	%pdfborder = {0 0 0},
	bookmarks=true,         % show bookmarks bar?
	unicode=false,          % non-Latin characters in Acrobat’s bookmarks
	pdftitle={Erstellung von Irrbildern zur Überlistung einer Verkehrsschilder erkennenden KI}
	pdfauthor={Leonhard Applis, Peter Bauer, Andreas Porada und Florian Stöckl},     % author
	pdfsubject={IT-Projektarbeit},   % subject of the document
	pdfcreator={Leonhard Applis, Peter Bauer, Andreas Porada und Florian Stöckl},   % creator of the document
	pdfproducer={Leonhard Applis, Peter Bauer, Andreas Porada und Florian Stöckl}, % producer of the document
	colorlinks=true,       % false: boxed links; true: colored links
	linkcolor=blue,          % color of internal links (change box color with inkbordercolor)
	citecolor=CadetBlue,        % color of links to bibliography
	filecolor=magenta,      % color of file links
	urlcolor=cyan,           % color of external links
}

% -------------------------------------------------------------------------------------------
%                     Persönliche Daten
% -------------------------------------------------------------------------------------------


\newcommand{\titel}{Erstellung von Irrbildern zur Überlistung einer Verkehrsschilder erkennenden KI}
\newcommand{\shellcmd}[1]{\\\indent\indent\texttt{\footnotesize\# #1}\\}
\newcommand{\untertitel}{}
\newcommand{\arbeit}{IT-Projekt Bericht}
\newcommand{\studiengang}{Informatik}
\newcommand{\autor}{Leonhard Applis, Peter Bauer, Andreas Porada und Florian Stöckl}
\newcommand{\matrikelnr}{2086307}
\newcommand{\abgabe}{15.03.2019}
\newcommand{\betreuerth}{Prof. Dr. Gallwitz}
\newcommand{\jahr}{2018}			% für Angabe im Copyright-Vermerk der Titelseite


% -------------------------------------------------------------------------------------------
%                     Beginn des Dokumenteninhalts
% -------------------------------------------------------------------------------------------


\begin{document}
\setcounter{secnumdepth}{3}					% Nummerierungstiefe fürs Inhaltsverzeichnis
\setcounter{tocdepth}{3}
\sffamily									% für die Titelei serifenlose Schrift verwenden

% ------------------------------ Titelei -----------------------------------------------------

\thispagestyle{plain}
\begin{titlepage}
\enlargethispage{3.5cm}
\sffamily 								% Serifenlose Grundschrift für die Titelseite einstellen
\begin{minipage}{\textwidth}
	\vspace{-2cm}
	\noindent \includegraphics[scale=0.2]{Images/Logo_TH.png} \hfill
\end{minipage} 
\begin{center}

\huge{\textsc{\textbf{\titel}}}\\[1.5ex]
\Large{\textbf{\untertitel}}\\[5ex]
\LARGE{\textbf{\ausarbeitung}}\\[2ex]
\normalsize{~}\\[3ex]
\Large{Master-Studiengang \textit{\studiengang}}\\[1ex]
\normalsize{Technische Hochschule Nürnberg Georg Simon Ohm}\\[5ex]
von\\[1ex] \autor \\[12ex]

\begin{tabular}{ll}
	Abgabedatum:					& \quad \abgabeTH \\ 
	
	Gutachter der Hochschule: & \quad \betreuerth \\ 
	%[6ex]%formerly 5ex
	
\end{tabular} 

\end{center}

\end{titlepage}
%Starting Correct Spacing here...
\onehalfspacing 				% erzeugt die Titelseite
\pagenumbering{Roman}						% große, römische Seitenzahlen für Titelei
\addchap*{Eidesstattliche Erklärung}
Wir versichern hiermit, dass der \arbeit~ mit dem Thema

\emph{\titel}\\
selbständig verfasst wurde und keine anderen als die angegebenen Quellen und Hilfsmittel benutzt wurden. Die Arbeit wurde bisher keiner anderen Prüfungsbehörde vorgelegt und auch nicht veröffentlicht.


Wir versichern zudem, dass die eingereichte elektronische Fassung mit der gedruckten Fassung übereinstimmt.\\[10ex]

Nürnberg, den \today \\[4ex]


\rule[-0.2cm]{10cm}{0.5pt} \\

\textsc{\autor} \\[10ex]
 				% Einbinden der eidestattlichen Erklärung
\chapter*{Abstract} %*-Variante sorgt dafür, das Abstract nicht im Inhaltsverzeichnis auftaucht
This paper uses several scientific approaches to show how convolutional neural networks can be tricked into recognizing and correctly classifying road signs.
In the field of autonomous driving, neural networks have drastically increased the recognition rate, but they are still susceptible to errors in the face of deliberately generated false images. Even without information about the underlying architecture (so-called black box attacks), attacks by otherwise generated false images should be possible.
The presented methods \textit{Degeneration}, \textit{Saliency Maps} and \textit{Gradient Ascent} are successfully applied to generate false images for attacks on an unknown neural network with the help of an own neural network which serves as \textit{Substitute}.
An attack is considered "'successful"' if the image is recognized as a street sign with a confidence greater than 90\%, which a human observer would not recognize as such.
~\newline
~\newline
\begin{flushleft}
	\begin{tabular}{lp{11cm}}
		\textbf{title:} & Fooling a TrafficSign-AI \\
		\textbf{authors:}  & \autor \\
		
		%einkommentieren für TH Abgabe
		%\textbf{reviewer TH:} & \betreuerth \\
		%[6ex]%formerly 5ex
	\end{tabular} 
\end{flushleft}


\chapter*{Kurzfassung} 
Diese Arbeit zeigt anhand von mehreren wissenschaftlichen Ansätzen, wie Convolutional Neural Networks zur Erkennung und korrekten Klassifikation von Straßenschildern überlistet werden können.
Im Bereich des Autonomen Fahren wurden mit Neuronalen Netzen die Erkennungsrate drastisch gesteigert, allerdings sind diese immer noch fehleranfällig gegenüber gezielt erzeugten Irrbildern. Selbst ohne Informationen über die unterliegende Architektur (sog. Black-Box Angriffe), sollen Angriffe durch anderweitig erzeugte Irrbilder möglich sein.
Die vorgestellten Verfahren \textit{Degeneration}, \textit{Saliency Maps} und \textit{Gradient Ascent} werden erfolgreich angewendet, um mithilfe eines eigenen Neuronalen Netz welches als \textit{Substitute} dient, Irrbilder für Angriffe auf ein unbekanntes Neuronales Netz zu erzeugen.
Ein Angriff gilt als "'erfolgreich", wenn das Bild mit einer Konfidenz größer als 90\% als Straßenschild erkannt ist, welches ein menschlicher Betrachter nicht als solches erkennen würde.

~\newline
\begin{flushleft}
	\begin{tabular}{lp{11cm}}
		Titel:&  \titel \\ 
		Authoren:&  \autor \\
		%einkommentieren für TH Abgabe
		%Prüfer der Hochschule: &  \betreuerth \\ 
		%[6ex]%formerly 5ex	
	\end{tabular} 
\end{flushleft}   				% Einbinden des Abstracts

\tableofcontents							% Erzeugen des Inhalsverzeichnisses
\printnomenclature[2.0cm]					% Erzeugen des Abkürzungsverzeichnisses
\listoffigures 								% Erzeugen des Abbildungsverzeichnisses 
%\listoftables 								% Erzeugen des Tabellenverzeichnisses
\pagebreak
% ------------------ Graphic Extension ------------------------------------------------------

% --------------------------------------------------------------------------------------------
%                    Inhalt der Bachelorarbeit
%---------------------------------------------------------------------------------------------
\pagenumbering{arabic}						% arabische Seitenzahlen für den Hauptteil
\pagestyle{fancy}					
\rmfamily

%% alle Abkürzungen, die in der Arbeit verwendet werden
\chapter*{Abkürzungsverzeichnis}
\addcontentsline{toc}{chapter}{Abkürzungsverzeichnis}
\begin{singlespace}
	
\begin{acronym}[Studienarbeit]
    \acro{DBMS}{Database Management System}
	\acro{ETL}{Extract-Transform-Load}
	\acro{KI}{Künstliche Intelligenz}
	\acro{SQL}{Structured Query Language}
	\acro{NN}{Neuronales Netz}
	\acro{DNN}{Deep Neural-Network}
	\acro{CNN}{Convolutional Neural-Network}
	\acro{GTSRB}{German Traffic Sign Recognition Benchmark}
	\acro{PPM}{Portable Pixmap Image}
	\acro{PNG}{Portable Network Graphic}
\end{acronym}

\end{singlespace}
\chapter{Einleitung}
\label{cha:Einleitung}
\setlength{\epigraphwidth}{4in}

%\section{Motivation}
Der Traum von einem autonom gelenkten Automobil ist so alt wie das Automobil selbst \cite{maurer_autonomes_2015}. Fabian Dröger schreibt dazu in seinem Beitrag "`Das automatisierte Fahren im gesellschaftswissenschaftlichen und kulturwissenschaftlichen Kontext"' in dem Sammelwerk "`Autonomes Fahren"' von Markus Maurer et al., wie die zunehmende Anzahl an Verkehrstoten in den USA zu Beginn des 20. Jahrhunderts in Verbindung mit den technischen Errungenschaften in der frühen Flugzeug- und Radiotechnik den Wunsch nach einem selbstfahrenden Automobil aufkommen ließen. Die Vision war, dass ein Automobil ähnlich wie ein Flugzeug durch einen Autopiloten in der Spur gehalten und gesteuert werden könnte. Für die Ansteuerung der mechanischen Teile setzte man auf eine Fernsteuerung mit Funk, die zu dieser Zeit im Bereich der \emph{Radioguidance} erforscht wurde.

Der aktuelle Stand der Technik zeigt, dass sich die Umsetzung dieser Vision schwieriger gestaltet, als zunächst angenommen. Anstelle einer autonomen Steuerung finden sich in heutigen Automobilen verschiedene Techniken zur Erhöhung der Fahrsicherheit und des Komforts. Beispiele sind Spurhalteassistenten, automatische Abstandshalter oder Einparkhilfen. Diese Funktionen unterstützen einen menschlichen Fahrer, ermöglichen jedoch noch kein selbstständiges Fahren.

Es wird aber weiterhin an der Entwicklung eines autonomen Fahrzeugs geforscht, wie die Vergabe von Forschungsgeldern\cite{bmbf-internetredaktion_auto_nodate} und Berichte von Automobil"-her"-stel"-lern\cite{bmw_autonomes_nodate} und der Presse\cite{efler_autonomes_2018} zeigen. Die Forschung im Bereich \ac{KI} hat mittlerweile einen Stand erreicht, der für die Automatisierung des Autos genutzt werden kann. Ein besonderer Fokus liegt hierbei auf dem Erkennen von Bilder aus der Umwelt mittels einer \ac{KI}. Im Straßenverkehr ist besonders die Erkennung von Straßenschildern von Bedeutung.


\section{Problemstellung}
Bei der Untersuchung von \acp{KI}, welche auf die Erkennung gewisser Muster mit einem Datensatz trainiert wurden, fiel auf, dass auch Muster, die für den Menschen in keinem erkennbaren Zusammenhang mit dem Trainingsmuster stehen, von der \ac{KI} mit hoher Konfidenz anerkannt werden. Das gezielte Ausnutzen dieses Fehlers wird in der Forschung als \textit{Adversarial Attack} bezeichnet \cite{DBLP:journals/corr/HuangPGDA17}.

Der Erfolg von \textit{Adversarial Attacks} liegt in der gezielten Stimulation von Gewichten, die die \ac{KI} basierend auf ihrem Trainingssatz festgelegt hat, beziehungsweise die vom Trainer definiert wurden. Auf diese Weise ist es möglich ein gewünschtes Feedback eines \acs{NN}s zu erhalten, obwohl das Muster für einen Menschen kaum oder gar nicht mit den Mustern des Trainingssatzes in Verbindung gebracht werden kann. Zum Beispiel haben erzeugte Fragmente zur Täuschung einer \ac{KI}, die für die Erkennung von Bildinhalten trainiert ist, selten etwas mit einem \textit{echten} Bild zu tun. Sie wirken eher wie Rauschen oder moderne Kunst. 

~\newline Das Problem an den präparierten Mustern einer \textit{Adversarial Attack} ist, dass ein Mensch schwer erkennen kann, ob ein derartiger Angriff auf die von ihm genutzte \ac{KI} unternommen wird. Er ordnet das manipulierte Muster nicht den relevanten Mustern zu, die seine \ac{KI} anerkennt.

Durch den steigenden Einsatz von Machine Learning in verschiedenen sensiblen Sektoren des täglichen Lebens, wie selbstfahrenden Autos, Terrorismusbekämpfung oder Betrugserkennung können Angriffe verheerende Schäden erzeugen und stellen ein lohnendes Ziel dar. Vor allem im Bereich des autonomen Fahrens, welcher aktuell geprägt ist durch eine Debatte über das \textit{Vertrauen in Technik} \cite{VertrauenTechnik}, können erfolgreiche Angriffe zu einem forschungsschädlichen Misstrauen führen und die Nutzung der Technik verhindern. 

~\newline Um gegen Adversarial Attacks vorzugehen, werden zunächst einige Bilder benötigt, die gezielt Gewichte einer \ac{KI} stimulieren und mit hoher Konfidenz anerkannt werden. Anschließend können, um das neuronale Netz zu härten, Tests durchgeführt werden und die Angriffe berücksichtigt werden. 

%Diese Irrbilder zu erzeugen stellt das Kernziel dieser Arbeit dar. 

\section{Ziel der Arbeit}
\label{sec:ZielDerArbeit}
Ziel dieser Arbeit ist es, Methoden und Herangehensweisen vorzustellen, um die Aufgabenstellung des Informaticup 2019 \todo{Hyperlink!} zu erfüllen: 

Hierbei soll ein neuronales Netz, welches sich hinter einer Webschnittstelle verbirgt und Verkehrsschilder erkennt, erfolgreich \textit{überlistet} werden - es sollen absichtlich Bilder erzeugt werden, welche für den Menschen keine Verkehrsschilder sind aber mit einer Konfidenz von über 90\% als solche erkannt werden. 

Die gefundenen Methoden sollen reproduzierbar sein und in einem Maße flexibel, um beliebig viele Irrbilder zu erzeugen. 

~\newline Der erweiterte Rahmen dieser Arbeit umfasst eine Dokumentation der Methoden sowie Verbesserungen und Schlussfolgerungen aus den Implementierungen zu ziehen. 

Ebenfalls geliefert werden alle Elemente, um die erzielten Ergebnisse zu reproduzieren und zu variieren. 

~\newline Nicht Ziel dieser Arbeit ist, einen Überblick über neuronale Netze, künstliche Intelligenz oder Bildbearbeitung zu vermitteln. 

~\newline Ebenfalls außerhalb dieser Arbeit liegt eine Auswertung, welche Bilder von einem Menschen als Verkehrsschilder erkannt werden. Die Aussagen über solche stützen sich ausschließlich auf die persönliche Einschätzung des Projektteams. 
\section{Aufbau der Arbeit}
Innerhalb dieser Arbeit werden zunächst in Kapitel \ref{cha:AnfAnalyse} Informationen über die Webschnittstelle gesammelt und aufbereitet. Die Webschnittstelle stellt eine zentrale Rahmenbedingung der vorliegenden Arbeit dar.


Im Abschnitt \ref{sec:EigenschaftenTrasi} wird hierfür zunächst der \ac{GTSRB} vorgestellt, welcher für das Training der Webschnittstelle verwendet wurde. Dieses Datenset bildet ebenfalls einen zentralen Ausgangspunkt für einige der verfolgten Ansätze.


Anschließend werden in Abschnitt \ref{sec:TrasiModell} die Eigenschaften des Modells zusammengefasst. 
Diese bestehen zum einen aus den offiziellen Angaben der Gesellschaft der Informatiker, zum anderen aus gewonnenen Erkenntnissen. Dieses Kapitel bildet die Grundlage, um die Schnittstelle einzuschätzen. 

~\newline Anschließend werden verschiedene Lösungsansätze vorgestellt, beginnend mit der \textit{Degeneration} in Kapitel \ref{cha:Degeneration}. 

Dieser Ansatz verändert iterativ ein Verkehrsschild, und behält die Änderungen bei, sollte der erzielte Score im akzeptablen Bereich liegen. Mit passender Bildveränderungen erzielen höhere Iterationen unkenntliche Ergebnisse. 

Innerhalb des Kapitels wird zunächst im Abschnitt \ref{sec:DegenerationKonzept} die Idee anhand von Pseudocode weiter erläutert und anschließend in Abschnitt  \ref{sec:DegenerationRemote} die Implementierung für die Webschnittstelle gezeigt. Die Ergebnisse liegen gesondert im Abschnitt \ref{sec:DegenerationErgebnisse} vor. 

Neben der Implementierung für die Webschnittstelle werden zum Abschluss des Kapitels in Abschnitt \ref{sec:DegenerationLokal} noch weitere Verbesserungen für eine lokale Implementierung vorgestellt, welche allerdings nicht für die Webschnittstelle tauglich waren.


~\newline Desweiteren werden in Kapitel \ref{cha:saliency} verschiedene Methoden zur Erzeugung von sog. \textit{Saliency Maps} (dt. Ausprägungskarte) vorgestellt. Unveränderte Bilder mit einer hohen Konfidenz werden ausgewählt, um die einzelnen Pixel hervorzuheben, welche für die Klassifikation den höchsten Einfluss hatten.

~\newline In Kapitel \ref{cha:gascent} wird das Verfahren des \textit{Gradient Ascent} beschrieben und evaluiert. Dabei wird mithilfe einer \textit{targeted Backpropagation} für jede Klasse ein Bild erzeugt, der die Funktion in Richtung der Zielklasse maximiert, bis die enthaltenen Merkmale eine hohe Konfidenz in der gewünschten Klasse ermöglichen.

~\newline Abschluss dieser Arbeit bildet im Kapitel \ref{cha:Schluss} ein Fazit über die gefundenen Methoden, sowie ein Ausblick auf weiterführende Arbeiten. 

\chapter{Ansatz 1 : Greyboxing}
\label{Cha:GreyBoxing}
\section{Konzept}

\section{Implementierung und erste Ergebnisse}

\section{Tuning}

\subsection{Iteration 2: Bessere Trainingsdaten Multiclass}

\subsection{Iteration 3: Verbesserte Bilderstellung}

\subsection{Iteration 4: Automatisierte Parameterisierung der Bilder}

\section{Endergebnisse}
\chapter{Ansatz 2 : Genetische Algorithmen [AnPe]}
\label{GenAlgo}
\chapter{Fazit}
\label{cha:Fazit} \label{cha:Schluss}
Nach einer abschließenden Zusammenfassung, werden in diesem Kapitel zusätzlich die verschiedenen Ansätze verglichen und bewertet. Aus dieser Diskussion werden die Einschränkungen und Möglichkeiten für weitere Arbeiten angesprochen.

\section{Zusammenfassung} ~\newline 
\todo{Ergebnisse Degen}

Desweiteren konnten Erfolge mit der Erzeugung von Saliency Maps in den geglätteten Varianten verbucht werden. Somit konnte bestätigt werden, dass die Visualisierung der Relevanten Pixel für ein Bild mit hoher Konfidenz geeignet sind, um als "minimale Beispiele" für das NN verwendet werden können. Die entsprechenden Verfahren erzeugten Täuschungen mit Konfidenzen >0.9, allerdings lassen sich keine Aussagen über die allgemeine Verlässlichkeit und Zielgerichtetheit treffen.

Zuletzt konnte auch mit dem \textit{Gradient Ascent} Verfahren sehr gute ergebnisse erzielt werden. für 10 von 43 Klassen konnten Täuschungen mit hohen Konfidenzen am Remote-NN erzielt werden. Es wird vermutet dass die Ausbeute mit weiterer Optimierung vergrößert werden kann oder auch echte Bilder für die Erzeugung der Täuschungen verwendet werden können.



\section{Diskussion}
Die Herausforderungen des Wettbewerbs wirken sich auf die Erfolge der Verfahren aus. Es wird vermutet das beide Verfahren \textit{Saliency Map} und \textit{Gradient Ascent} bessere ergebnisse liefern Könnten, wenn das verwendete Bild größer als $64\times64$ wäre. Desweiteren kann nichts über die Validierungsgenauigkeit des Trasi-NN gesagt werden, weshalb auch die nicht zielgerichteten Täuschungsbilder als gutes Ergebnis betitelt werden. 


Im Vergleich der Methoden \textit{Saliency Map} und \textit{Gradient Ascent}, kann das letztere bevorzugt werden. 



hier kommen rein
vergleich der ergebnisse (reproduzierbarkeit, gezielt/random, qualität der bilder)
- einschränkungen probleme unserer ansätze, 


wichtig: eigene fehler oder einschränkungen der methodik erkennen.



\section{Weiterführende Arbeiten}~\newline 
Die Ergebnisse dieser Arbeit liefern weitere Ansätze für zukünftige Aufgaben. Zum einen können die verwendeten Ansätze individuell weiter Optimiert werden, bezüglich des selbsterstellten lokalen Neuronalem Netz und der Algorithmik bzw. deren Parameter. Desweiteren können verfahren entwickelt werden, die sich aller Methoden gezielt bedienen. 
%content
%Vergleich mit einem Framework, wie zum Beispiel CleverHans \footnote, Vergleich von einem gut und schlecht trainierten Modell, wie arg unterscheiden sich daraus generierte Bilder (lohnt es sich), wie wichtig ist die Bildgröße, 
\newpage\newpage

% ---------------------------- Literaturverzeichnis ----------------------------------------------

\bibliographystyle{plain}
\bibliography{src}


% ------------------------------- Anhang ---------------------------------------------------------

\begin{appendix}
\clearpage
\pagenumbering{Roman}						% römische Seitenzahlen für Anhang
\end{appendix}


\end{document}
